% Public question bank LTL specification appendix.
% Generated by tools/render_spec_tex.py from codesys_service/specs/*_spec.md.
% Do not edit by hand: edit the markdown and re-run the script.
% Fragment: no preamble. Needs \usepackage{longtable,booktabs,array}.

\section{Public Question Bank LTL Specifications}
\label{app:public_ltl}

This appendix gives the formal specification of the 20 public questions. Each task has a set of properties written in linear temporal logic over the interface variables of its function block, one state per scan cycle: the inputs read in that cycle and the outputs computed in it. The 20 tasks carry 264 properties in total.

The operators are $\mathbf{G}$ (always), $\mathbf{F}$ (eventually), $\mathbf{X}$ (next cycle), $\wedge$ (and), $\vee$ (or), $\neg$ (not) and $\rightarrow$ (implies). Formulas are printed in Structured Text spelling with \textbf{AND}, \textbf{OR} and \textbf{NOT}, because the atoms are the ST identifiers of the interface and are kept identical to the source files. The abbreviation

\begin{center}
\texttt{rise(b) := NOT b AND X b}
\end{center}

\noindent is used throughout for \emph{b goes FALSE to TRUE across the step into the next cycle}. Operator buttons, start buttons and completion signals are edge-triggered and always appear under \texttt{rise}; sensors report a level and never do. Durations (\texttt{DebounceTime}, \texttt{CycleTime}) are accumulated time and are not edge-detected. Outputs are evaluated in the same scan cycle as the transition that causes them: a step taken in cycle $N$ is already visible in the outputs of cycle $N$. Where a formula needs the value in the following cycle it says so with an explicit $\mathbf{X}$.

Each task section also lists the assumptions taken where the natural-language requirement is silent or ambiguous, and the validation result of the accompanying test suite. Test cases are in the unified format of Section~\ref{app:caseformat}; property-to-case coverage is given in the per-task \texttt{.testcase.json} files.

{\small
\begin{longtable}{@{}lrlr@{}}
\toprule
\textbf{Task} & \textbf{Properties} & \textbf{Test cases} & \textbf{Flagged ambiguities} \\
\midrule
\endhead
\texttt{E01\_Filling} & 15 & 6 & 4 \\
\texttt{E02\_DualEdge} & 12 & 4 & 2 \\
\texttt{E03\_Hysteresis} & 12 & 5 & 5 \\
\texttt{E04\_JamPump} & 11 & 4 & 4 \\
\texttt{E05\_WeldGap} & 10 & 3 & 3 \\
\texttt{E06\_Blower} & 9 & 3 & 3 \\
\texttt{H01\_Proportion} & 22 & 6 & 6 \\
\texttt{H02\_FlexProcess} & 11 & 3 & 5 \\
\texttt{H03\_RecipeMgr} & 13 & 3 & 5 \\
\texttt{H04\_FIFO} & 12 & 5 & 5 \\
\texttt{H05\_MechArm} & 19 & 7 & 6 \\
\texttt{H06\_Production} & 21 & 6 & 6 \\
\texttt{M01\_PumpSwitch} & 8 & 3 & 4 \\
\texttt{M02\_AnalogBatch} & 12 & 3 & 4 \\
\texttt{M03\_MovingAvg} & 11 & 5 & 4 \\
\texttt{M04\_TempCtrl} & 12 & 4 & 5 \\
\texttt{M05\_ProdLine} & 18 & 6 & 5 \\
\texttt{M06\_SlabThick} & 14 & 5 & 4 \\
\texttt{M07\_HVAC} & 13 & 3 & 3 \\
\texttt{M08\_IPv4} & 9 & 4 & 4 \\
\midrule
\textbf{Total} & \textbf{264} & \textbf{88} & \\
\bottomrule
\end{longtable}
}

\subsection{E01\_Filling --- Automated Filling Line Control}
\label{spec:E01_Filling}

\noindent \texttt{E01\_Filling}; 15 properties; source \texttt{specs/E01\_Filling\_spec.md}.

{\small
\begin{longtable}{@{}p{0.05\linewidth}p{0.13\linewidth}p{0.30\linewidth}>{\raggedright\arraybackslash}p{0.42\linewidth}@{}}
\toprule
\textbf{ID} & \textbf{Category} & \textbf{Requirement} & \textbf{LTL} \\
\midrule
\endhead
P1 & Init & At power-up all outputs are OFF. & \texttt{idle} (holds in the initial state) \\
P2 & Mutual exclusion & At most one stage output is active at any time. & \texttt{G(idle \textbf{OR} onlyPump \textbf{OR} onlyValve \textbf{OR} onlyCapping \textbf{OR} onlyPacking \textbf{OR} onlyLight)} \\
P3 & Start & In idle, a bottle at the sensor starts the cleaning pump. & \texttt{G((idle \textbf{AND} X rise(bottleSensor)) -> X Pump\_Motor)} \\
P4 & Transition & During cleaning, the cleaning confirm button stops the pump and opens the filling valve. & \texttt{G((Pump\_Motor \textbf{AND} X rise(cleaningConfirmButton)) -> X(Filling\_Valve \textbf{AND} \textbf{NOT} Pump\_Motor))} \\
P5 & Transition & During filling, the filling confirm button closes the valve and starts capping. & \texttt{G((Filling\_Valve \textbf{AND} X rise(fillingConfirmButton)) -> X(Capping\_Machine \textbf{AND} \textbf{NOT} Filling\_Valve))} \\
P6 & Transition & During capping, the capping confirm button stops capping and starts packing. & \texttt{G((Capping\_Machine \textbf{AND} X rise(cappingConfirmButton)) -> X(Packing\_Machine \textbf{AND} \textbf{NOT} Capping\_Machine))} \\
P7 & Transition & During packing, the packing confirm button stops packing and lights the completion light. & \texttt{G((Packing\_Machine \textbf{AND} X rise(packingConfirmButton)) -> X(Completion\_Light \textbf{AND} \textbf{NOT} Packing\_Machine))} \\
P8 & Reset & When the light is on, finishedButton turns it off and returns to idle. & \texttt{G((Completion\_Light \textbf{AND} X rise(finishedButton)) -> X idle)} \\
P9 & Latching & The pump keeps running until the cleaning confirm button is pressed (even if the sensor is released). & \texttt{G((Pump\_Motor \textbf{AND} X \textbf{NOT} cleaningConfirmButton) -> X Pump\_Motor)} \\
P10 & Latching & The valve stays open until the filling confirm button is pressed. & \texttt{G((Filling\_Valve \textbf{AND} X \textbf{NOT} fillingConfirmButton) -> X Filling\_Valve)} \\
P11 & Latching & Capping stays active until the capping confirm button is pressed. & \texttt{G((Capping\_Machine \textbf{AND} X \textbf{NOT} cappingConfirmButton) -> X Capping\_Machine)} \\
P12 & Latching & Packing stays active until the packing confirm button is pressed. & \texttt{G((Packing\_Machine \textbf{AND} X \textbf{NOT} packingConfirmButton) -> X Packing\_Machine)} \\
P13 & Latching & The light stays on until finishedButton is pressed. & \texttt{G((Completion\_Light \textbf{AND} X \textbf{NOT} finishedButton) -> X Completion\_Light)} \\
P14 & No spurious start & In idle, nothing starts without a bottle. & \texttt{G((idle \textbf{AND} X \textbf{NOT} bottleSensor) -> X idle)} \\
P15 & Sequence integrity & Each stage can only be entered from its predecessor (e.g. the valve never opens unless the pump was running in the previous cycle). & \texttt{G(X Filling\_Valve -> (Pump\_Motor \textbf{OR} Filling\_Valve))}, and analogously for Capping\_Machine, Packing\_Machine, Completion\_Light \\
\bottomrule
\end{longtable}
}

\noindent\textbf{Assumptions and requirement gaps.}
{\small\begin{enumerate}
\item \textbf{Level vs. edge.}  The requirement only says the operator "triggers" the signal. The team adopted \textbf{rising-edge} semantics for the confirmation buttons; the test cases therefore release the button between operations. \emph{(flagged as ambiguous)}
\item \textbf{Out-of-sequence buttons} are not mentioned; we assume they are ignored (P9--P14), which is the only interpretation consistent with "each step triggered by its own signal".
\item \textbf{A new bottle while the light is on} is not specified; we assume the cycle cannot restart until finishedButton (P13), since the operator must pick up the finished bottle first.
\item "Continues to work" (pump) is formalised as latching (P9).
\end{enumerate}}

\noindent\textbf{Validation.} reference implementation, simulation 6/6 cases PASS; mutation testing 2/2 mutants detected; 6 cases / 57 scan cycles.

\subsection{E02\_DualEdge --- Dual-Word Edge Detection}
\label{spec:E02_DualEdge}

\noindent \texttt{E02\_DualEdge}; 12 properties; source \texttt{specs/E02\_DualEdge\_spec.md}.

{\small
\begin{longtable}{@{}p{0.05\linewidth}p{0.13\linewidth}p{0.30\linewidth}>{\raggedright\arraybackslash}p{0.42\linewidth}@{}}
\toprule
\textbf{ID} & \textbf{Category} & \textbf{Requirement} & \textbf{LTL} \\
\midrule
\endhead
P1 & Rising mask & A bit of \texttt{risingBits} is set exactly when that bit went 0$\rightarrow$1. & \texttt{G(X risingBits = (X value \textbf{AND} \textbf{NOT} value))} \\
P2 & Falling mask & A bit of \texttt{fallingBits} is set exactly when that bit went 1$\rightarrow$0. & \texttt{G(X fallingBits = (value \textbf{AND} \textbf{NOT} X value))} \\
P3 & Change flag & \texttt{hasChanged} is TRUE exactly when the value differs from the previous cycle. & \texttt{G(X hasChanged <-> (X value <> value))} \\
P4 & Rising flag & \texttt{hasRisingEdges} is TRUE exactly when at least one rising edge occurred. & \texttt{G(hasRisingEdges <-> (risingBits <> 0))} \\
P5 & Falling flag & \texttt{hasFallingEdges} is TRUE exactly when at least one falling edge occurred. & \texttt{G(hasFallingEdges <-> (fallingBits <> 0))} \\
P6 & Rising count & \texttt{noOfRisingBits} equals the number of set bits in \texttt{risingBits}. & \texttt{G(noOfRisingBits = popcount(risingBits))} \\
P7 & Falling count & \texttt{noOfFallingBits} equals the number of set bits in \texttt{fallingBits}. & \texttt{G(noOfFallingBits = popcount(fallingBits))} \\
P8 & Disjointness & A bit cannot rise and fall in the same cycle. & \texttt{G((risingBits \textbf{AND} fallingBits) = 0)} \\
P9 & Consistency & \texttt{hasChanged} holds iff at least one edge of either direction occurred. & \texttt{G(hasChanged <-> (hasRisingEdges \textbf{OR} hasFallingEdges))} \\
P10 & No change & If the value is unchanged, all edge outputs are zero. & \texttt{G(\textbf{NOT} hasChanged -> (risingBits = 0 \textbf{AND} fallingBits = 0 \textbf{AND} noOfRisingBits = 0 \textbf{AND} noOfFallingBits = 0))} \\
P11 & Count bounds & Both counts are at most 32. & \texttt{G(noOfRisingBits <= 32 \textbf{AND} noOfFallingBits <= 32)} \\
P12 & First invocation  & On the first cycle the previous value is treated as 0, so a non-zero input produces rising edges only. & \texttt{(value <> 0) -> (hasRisingEdges \textbf{AND} \textbf{NOT} hasFallingEdges)} (initial state) \\
\bottomrule
\end{longtable}
}

\noindent\textbf{Assumptions and requirement gaps.}
{\small\begin{enumerate}
\item \textbf{P12}: the requirement does not define the state before the first invocation; we assume the stored previous value is 0 (standard IEC initialisation).
\item Outputs are computed every cycle; there is no enable input.
\end{enumerate}}

\noindent\textbf{Validation.} reference implementation, simulation 4/4 cases PASS; mutation testing 2/2 mutants detected; 4 cases / 18 scan cycles.

\subsection{E03\_Hysteresis --- Alarm Handling}
\label{spec:E03_Hysteresis}

\noindent \texttt{E03\_Hysteresis}; 12 properties; source \texttt{specs/E03\_Hysteresis\_spec.md}.

{\small
\begin{longtable}{@{}p{0.05\linewidth}p{0.13\linewidth}p{0.30\linewidth}>{\raggedright\arraybackslash}p{0.42\linewidth}@{}}
\toprule
\textbf{ID} & \textbf{Category} & \textbf{Requirement} & \textbf{LTL} \\
\midrule
\endhead
P1 & Error handling & If loLevel $\geq$ hiLevel, error is TRUE and status is 16\#8001. & \texttt{G((loLevel >= hiLevel) -> (error \textbf{AND} status = 16\#8001))} \\
P2 & Error handling & If the limits are ordered but hystWindow $\geq$ hiLevel $-$ loLevel, error is TRUE and status is 16\#8002. & \texttt{G((loLevel < hiLevel \textbf{AND} hystWindow >= hiLevel - loLevel) -> (error \textbf{AND} status = 16\#8002))} \\
P3 & Alarm trigger & With valid parameters, Value > hiLevel sets hiAlarm. & \texttt{G((valid \textbf{AND} Value > hiLevel) -> hiAlarm)} \\
P4 & Alarm release & With valid parameters, Value < hiLevel $-$ hystWindow clears hiAlarm. & \texttt{G((valid \textbf{AND} Value < hiLevel - hystWindow) -> \textbf{NOT} hiAlarm)} \\
P5 & Alarm trigger & With valid parameters, Value < loLevel sets loAlarm. & \texttt{G((valid \textbf{AND} Value < loLevel) -> loAlarm)} \\
P6 & Alarm release & With valid parameters, Value > loLevel + hystWindow clears loAlarm. & \texttt{G((valid \textbf{AND} Value > loLevel + hystWindow) -> \textbf{NOT} loAlarm)} \\
P7 & Hysteresis (hold) & An active high alarm stays active while Value remains $\geq$ hiLevel $-$ hystWindow. & \texttt{G((hiAlarm \textbf{AND} X(valid \textbf{AND} Value >= hiLevel - hystWindow)) -> X hiAlarm)} \\
P8 & Hysteresis (hold) & An active low alarm stays active while Value remains $\leq$ loLevel + hystWindow. & \texttt{G((loAlarm \textbf{AND} X(valid \textbf{AND} Value <= loLevel + hystWindow)) -> X loAlarm)} \\
P9 & No false alarm & An inactive high alarm is not set while Value $\leq$ hiLevel. & \texttt{G((\textbf{NOT} hiAlarm \textbf{AND} X(valid \textbf{AND} Value <= hiLevel)) -> X \textbf{NOT} hiAlarm)} \\
P10 & No false alarm & An inactive low alarm is not set while Value $\geq$ loLevel. & \texttt{G((\textbf{NOT} loAlarm \textbf{AND} X(valid \textbf{AND} Value >= loLevel)) -> X \textbf{NOT} loAlarm)} \\
P11 & Normal status  & With valid parameters, no error is reported. & \texttt{G(valid -> (\textbf{NOT} error \textbf{AND} status = 16\#0000))} \\
P12 & Mutual exclusion & With valid parameters the two alarms are never active together. & \texttt{G(valid -> \textbf{NOT} (hiAlarm \textbf{AND} loAlarm))} \\
\bottomrule
\end{longtable}
}

\noindent\textbf{Assumptions and requirement gaps.}
{\small\begin{enumerate}
\item \textbf{P11 is an assumption.} The requirement does not specify \texttt{error}/\texttt{status} when the parameters are valid; we assume \texttt{error = FALSE} and \texttt{status = 16\#0000} (common Siemens/CODESYS convention).
\item \textbf{Error priority.} When both error conditions hold, the requirement says the program exits on 16\#8001, so 16\#8001 takes priority (P2 is guarded by \texttt{loLevel < hiLevel}).
\item \textbf{Alarm outputs during an error} are not specified ("exit the program"), so no property constrains them.
\item \textbf{Boundary values.} Triggering uses strict comparisons (\texttt{>} / \texttt{<}), so \texttt{Value = hiLevel} does not trigger and \texttt{Value = hiLevel $-$ hystWindow} does not release, exactly as worded.
\item \textbf{P12} follows from P3--P10 with valid parameters (hiLevel $-$ hystWindow > loLevel), and is included as an explicit safety check.
\end{enumerate}}

\noindent\textbf{Validation.} reference implementation, simulation 5/5 cases PASS; mutation testing 2/2 mutants detected; 5 cases / 23 scan cycles.

\subsection{E04\_JamPump --- Jam Pump Debouncing and Runtime Maintenance}
\label{spec:E04_JamPump}

\noindent \texttt{E04\_JamPump}; 11 properties; source \texttt{specs/E04\_JamPump\_spec.md}.

{\small
\begin{longtable}{@{}p{0.05\linewidth}p{0.13\linewidth}p{0.30\linewidth}>{\raggedright\arraybackslash}p{0.42\linewidth}@{}}
\toprule
\textbf{ID} & \textbf{Category} & \textbf{Requirement} & \textbf{LTL} \\
\midrule
\endhead
P1 & Safety & An emergency stop immediately disables the run output. & \texttt{G(\textbf{NOT} PASTE\_E\_STOP\_NC -> \textbf{NOT} VFD\_RUN)} \\
P2 & Safety & The pump never runs while the viscosity signal is FALSE. & \texttt{G(\textbf{NOT} VISCOSITY\_OK -> \textbf{NOT} VFD\_RUN)} \\
P3 & Debounce & The pump runs only after the viscosity condition has been continuously valid for at least \texttt{DebounceTime}. & \texttt{G(VFD\_RUN -> debAcc >= DebounceTime)} \\
P4 & Debounce & When the viscosity signal or the e-stop goes false, the accumulated debounce time is cleared. & \texttt{G((\textbf{NOT} VISCOSITY\_OK \textbf{OR} \textbf{NOT} PASTE\_E\_STOP\_NC) -> X(debAcc = 0 \textbf{OR} debAcc = CycleTime))} \\
P5 & Liveness & While both conditions hold continuously, the pump starts once \texttt{DebounceTime} has elapsed. & \texttt{G((VISCOSITY\_OK \textbf{AND} PASTE\_E\_STOP\_NC \textbf{AND} debAcc >= DebounceTime) -> VFD\_RUN)} \\
P6 & Runtime & Runtime accumulates exactly by \texttt{CycleTime} while the pump runs. & \texttt{G(VFD\_RUN -> X PUMP\_RUNTIME\_ACC = PUMP\_RUNTIME\_ACC + CycleTime)} \\
P7 & Runtime & Runtime is preserved while the pump is stopped. & \texttt{G(\textbf{NOT} VFD\_RUN -> X PUMP\_RUNTIME\_ACC = PUMP\_RUNTIME\_ACC)} \\
P8 & Maintenance & The maintenance indication is set when the accumulated runtime reaches 500 hours. & \texttt{G(PUMP\_RUNTIME\_ACC >= 500h -> PUMP\_MAINT\_DUE)} \\
P9 & Maintenance & The indication is never set below 500 hours. & \texttt{G(PUMP\_RUNTIME\_ACC < 500h -> \textbf{NOT} PUMP\_MAINT\_DUE)} \\
P10 & Independence & The maintenance indication never affects the run output. & \texttt{G(PUMP\_MAINT\_DUE \textbf{AND} VISCOSITY\_OK \textbf{AND} PASTE\_E\_STOP\_NC \textbf{AND} debAcc >= DebounceTime -> VFD\_RUN)} \\
P11 & Init & At power-up the run output is OFF, runtime zero and the indication FALSE. & \texttt{\textbf{NOT} VFD\_RUN \textbf{AND} PUMP\_RUNTIME\_ACC = 0 \textbf{AND} \textbf{NOT} PUMP\_MAINT\_DUE} (initial state) \\
\bottomrule
\end{longtable}
}

\noindent\textbf{Assumptions and requirement gaps.}
{\small\begin{enumerate}
\item \textbf{Units}: \texttt{TIME} inputs are in milliseconds. \texttt{PUMP\_RUNTIME\_ACC} is compared in \textbf{milliseconds}, matching \texttt{TIME\_TO\_LTIME} in the project toolchain (IEC nominally defines LTIME in nanoseconds).
\item \textbf{Debounce accumulation}  The requirement supplies \texttt{CycleTime} as "the actual elapsed interval", which allows two readings: accumulate \texttt{CycleTime} per invocation (simulated time), or use a \texttt{TON} timer (wall-clock). The team adopted the \textbf{TON} reading, so the test cases drive the virtual scan clock through \texttt{hold\_ms}. \emph{(flagged as ambiguous)}
\item \textbf{Maintenance flag latching}: the requirement does not mention a reset; we assume it stays TRUE once set.  \emph{(flagged as ambiguous)}
\item \textbf{500 hours} = 1 800 000 000 000 000 ns.
\end{enumerate}}

\noindent\textbf{Validation.} reference implementation, simulation 4/4 cases PASS; mutation testing 2/2 mutants detected; 4 cases / 20 scan cycles.

\subsection{E05\_WeldGap --- Weld Gap Alarm and Startup Inhibition}
\label{spec:E05_WeldGap}

\noindent \texttt{E05\_WeldGap}; 10 properties; source \texttt{specs/E05\_WeldGap\_spec.md}.

{\small
\begin{longtable}{@{}p{0.05\linewidth}p{0.13\linewidth}p{0.30\linewidth}>{\raggedright\arraybackslash}p{0.42\linewidth}@{}}
\toprule
\textbf{ID} & \textbf{Category} & \textbf{Requirement} & \textbf{LTL} \\
\midrule
\endhead
P1 & Safety & The emergency stop immediately clears the alarm and has priority over everything else. & \texttt{G(\textbf{NOT} Weld\_E\_Stop\_NC -> \textbf{NOT} Operator\_Alert)} \\
P2 & Inhibition & During the first 5 s after power-up the alarm is suppressed. & \texttt{G(initElapsed < 5s -> \textbf{NOT} Operator\_Alert)} \\
P3 & Inhibition flag & \texttt{Weld\_Soft\_Start\_Active} is TRUE exactly during the power-up inhibition period. & \texttt{G(Weld\_Soft\_Start\_Active <-> (initElapsed < 5s \textbf{AND} Weld\_E\_Stop\_NC))} \\
P4 & Inhibition & Debounce time does not accumulate during inhibition. & \texttt{G(initElapsed < 5s -> debAcc = 0)} \\
P5 & Debounce & The alarm is raised only after \texttt{Gap\_Exceeded} has been continuously TRUE for \texttt{DebounceTime}. & \texttt{G(Operator\_Alert -> debAcc >= DebounceTime)} \\
P6 & Debounce & A FALSE gap signal switches the alarm off and clears the timing. & \texttt{G(\textbf{NOT} Gap\_Exceeded -> (\textbf{NOT} Operator\_Alert \textbf{AND} X debAcc <= CycleTime))} \\
P7 & Liveness & After initialisation, with the e-stop inactive and the gap signal held for \texttt{DebounceTime}, the alarm is raised. & \texttt{G((Weld\_E\_Stop\_NC \textbf{AND} initElapsed >= 5s \textbf{AND} Gap\_Exceeded \textbf{AND} debAcc >= DebounceTime) -> Operator\_Alert)} \\
P8 & Priority & The emergency stop also clears incomplete debounce timing. & \texttt{G(\textbf{NOT} Weld\_E\_Stop\_NC -> X debAcc <= CycleTime)} \\
P9 & Init & At power-up the alarm is FALSE and the inhibition flag TRUE. & \texttt{\textbf{NOT} Operator\_Alert \textbf{AND} Weld\_Soft\_Start\_Active} (initial state) \\
P10 & Scope & \texttt{Weld\_Soft\_Start\_Active} is FALSE after the inhibition period and never returns. & \texttt{G(initElapsed >= 5s -> \textbf{NOT} Weld\_Soft\_Start\_Active)} \\
\bottomrule
\end{longtable}
}

\noindent\textbf{Assumptions and requirement gaps.}
{\small\begin{enumerate}
\item \textbf{Inhibition timing during an emergency stop}: the requirement says the e-stop has priority over initialisation. We assume the 5 s inhibition timer does not advance while the e-stop is active, and \texttt{Weld\_Soft\_Start\_Active} is FALSE then (P3).  \emph{(flagged as ambiguous)}
\item \textbf{Alarm release} is immediate when the gap clears ("switch off the alarm"), with no minimum hold.
\item \textbf{Inhibition is one-shot}: it applies only after power-up, not after an e-stop release (P10).
\end{enumerate}}

\noindent\textbf{Validation.} reference implementation, simulation 3/3 cases PASS; mutation testing 2/2 mutants detected; 3 cases / 18 scan cycles.

\subsection{E06\_Blower --- Aeration Blower Delayed Start}
\label{spec:E06_Blower}

\noindent \texttt{E06\_Blower}; 9 properties; source \texttt{specs/E06\_Blower\_spec.md}.

{\small
\begin{longtable}{@{}p{0.05\linewidth}p{0.13\linewidth}p{0.30\linewidth}>{\raggedright\arraybackslash}p{0.42\linewidth}@{}}
\toprule
\textbf{ID} & \textbf{Category} & \textbf{Requirement} & \textbf{LTL} \\
\midrule
\endhead
P1 & Safety & The blower never runs during an emergency stop. & \texttt{G(\textbf{NOT} AERATION\_E\_STOP\_NC -> \textbf{NOT} BLOWER\_RUN)} \\
P2 & Safety & The blower never runs without a low-dissolved-oxygen demand. & \texttt{G(\textbf{NOT} DO\_LOW -> \textbf{NOT} BLOWER\_RUN)} \\
P3 & Delay & The blower runs only after both conditions have held continuously for 5 s. & \texttt{G(BLOWER\_RUN -> delayAcc >= 5s)} \\
P4 & Reset & Loss of demand or an emergency stop resets the start timer. & \texttt{G((\textbf{NOT} DO\_LOW \textbf{OR} \textbf{NOT} AERATION\_E\_STOP\_NC) -> X delayAcc <= CycleTime)} \\
P5 & Liveness & While both conditions hold for 5 s the blower starts. & \texttt{G((DO\_LOW \textbf{AND} AERATION\_E\_STOP\_NC \textbf{AND} delayAcc >= 5s) -> BLOWER\_RUN)} \\
P6 & Counting & The start counter increments exactly on each FALSE$\rightarrow$TRUE transition of \texttt{BLOWER\_RUN}. & \texttt{G((\textbf{NOT} BLOWER\_RUN \textbf{AND} X BLOWER\_RUN) -> X AERATION\_BLOWER\_CYCLE\_COUNT = AERATION\_BLOWER\_CYCLE\_COUNT + 1)} \\
P7 & Counting & The counter does not change without such a transition. & \texttt{G(\textbf{NOT} (\textbf{NOT} BLOWER\_RUN \textbf{AND} X BLOWER\_RUN) -> X AERATION\_BLOWER\_CYCLE\_COUNT = AERATION\_BLOWER\_CYCLE\_COUNT)} \\
P8 & Init & At power-up the blower is OFF and the count is zero. & \texttt{\textbf{NOT} BLOWER\_RUN \textbf{AND} AERATION\_BLOWER\_CYCLE\_COUNT = 0} (initial state) \\
P9 & Monotonic & The counter never decreases. & \texttt{G(X AERATION\_BLOWER\_CYCLE\_COUNT >= AERATION\_BLOWER\_CYCLE\_COUNT)} \\
\bottomrule
\end{longtable}
}

\noindent\textbf{Assumptions and requirement gaps.}
{\small\begin{enumerate}
\item \textbf{Full delay after any interruption}: explicitly required ("a new demand must wait the full 5 seconds again").
\item \textbf{Counter overflow} of \texttt{UINT} is not specified and not tested.
\item \textbf{Emergency stop does not affect the counter} (only actual starts are counted).
\end{enumerate}}

\noindent\textbf{Validation.} reference implementation, simulation 3/3 cases PASS; mutation testing 2/2 mutants detected; 3 cases / 19 scan cycles.

\subsection{H01\_Proportion --- Material Proportioning Process Control}
\label{spec:H01_Proportion}

\noindent \texttt{H01\_Proportion}; 22 properties; source \texttt{specs/H01\_Proportion\_spec.md}.

{\small
\begin{longtable}{@{}p{0.05\linewidth}p{0.13\linewidth}p{0.30\linewidth}>{\raggedright\arraybackslash}p{0.42\linewidth}@{}}
\toprule
\textbf{ID} & \textbf{Category} & \textbf{Requirement} & \textbf{LTL} \\
\midrule
\endhead
P1 & Safety & Whenever the emergency stop is triggered, every device is OFF in the same cycle. & \texttt{G(\textbf{NOT} estop -> allOff)} \\
P2 & Init & All devices are OFF at power-up. & \texttt{allOff} (initial state) \\
P3 & No auto-resume & After an emergency stop is released, nothing restarts by itself; in automatic mode a new autoStart is required. & \texttt{G((\textbf{NOT} estop \textbf{AND} X(auto \textbf{AND} \textbf{NOT} autoStart)) -> X allOff)} \\
P4 & Toggle ON & A rising edge of valveAStart while valve A is OFF turns it ON (same for B, C, D, mixer). & \texttt{G((manual \textbf{AND} \textbf{NOT} valveARun \textbf{AND} rise(valveAStart) \textbf{AND} X manual) -> X valveARun)} \\
P5 & Toggle OFF & A rising edge of valveAStart while valve A is ON turns it OFF (same for B, C, D, mixer). & \texttt{G((manual \textbf{AND} valveARun \textbf{AND} rise(valveAStart) \textbf{AND} X manual) -> X \textbf{NOT} valveARun)} \\
P6 & Hold & Without a rising edge the state of valve A is unchanged (holding the button does not re-toggle). & \texttt{G((manual \textbf{AND} X(manual \textbf{AND} \textbf{NOT} rise\_now(valveAStart))) -> (valveARun <-> X valveARun))} where \texttt{rise\_now} is the edge evaluated in the next cycle \\
P7 & Independence & Only the device whose button has a rising edge changes state. & follows from P4--P6 applied per device \\
P8 & Auto inputs ignored & autoStart, levelSensor and mixMotorComplete have no effect in manual mode. & implied by P6 (the only way to change an output in manual mode is a button edge) \\
P9 & Start & When the process is idle, autoStart opens valve A. & \texttt{G((auto \textbf{AND} allOff \textbf{AND} X(auto \textbf{AND} rise(autoStart))) -> X valveARun)} \\
P10 & Idle & When idle and autoStart is not pressed, nothing starts. & \texttt{G((auto \textbf{AND} allOff \textbf{AND} X(auto \textbf{AND} \textbf{NOT} autoStart)) -> X allOff)} \\
P11 & A $\rightarrow$ B & Valve A closes and valve B opens when the level reaches 50 cm. & \texttt{G((auto \textbf{AND} valveARun \textbf{AND} X(auto \textbf{AND} levelSensor >= 50)) -> X(\textbf{NOT} valveARun \textbf{AND} valveBRun))} \\
P12 & A hold & Valve A stays open below 50 cm. & \texttt{G((auto \textbf{AND} valveARun \textbf{AND} X(auto \textbf{AND} levelSensor < 50)) -> X valveARun)} \\
P13 & B $\rightarrow$ mixer (two materials) & With processMode FALSE, valve B closes and the mixer starts at 80 cm; valve C is not opened. & \texttt{G((auto \textbf{AND} valveBRun \textbf{AND} X(auto \textbf{AND} levelSensor >= 80 \textbf{AND} \textbf{NOT} processMode)) -> X(\textbf{NOT} valveBRun \textbf{AND} mixMotorRun \textbf{AND} \textbf{NOT} valveCRun))} \\
P14 & B $\rightarrow$ C (three materials) & With processMode TRUE, valve B closes and valve C opens at 80 cm; the mixer does not start yet. & \texttt{G((auto \textbf{AND} valveBRun \textbf{AND} X(auto \textbf{AND} levelSensor >= 80 \textbf{AND} processMode)) -> X(\textbf{NOT} valveBRun \textbf{AND} valveCRun \textbf{AND} \textbf{NOT} mixMotorRun))} \\
P15 & B hold & Valve B stays open below 80 cm. & \texttt{G((auto \textbf{AND} valveBRun \textbf{AND} X(auto \textbf{AND} levelSensor < 80)) -> X valveBRun)} \\
P16 & C $\rightarrow$ mixer & Valve C closes and the mixer starts at 100 cm. & \texttt{G((auto \textbf{AND} valveCRun \textbf{AND} X(auto \textbf{AND} levelSensor >= 100)) -> X(\textbf{NOT} valveCRun \textbf{AND} mixMotorRun))} \\
P17 & C hold & Valve C stays open below 100 cm. & \texttt{G((auto \textbf{AND} valveCRun \textbf{AND} X(auto \textbf{AND} levelSensor < 100)) -> X valveCRun)} \\
P18 & Mixer $\rightarrow$ D & When mixing is complete the mixer stops and valve D opens; the process ends. & \texttt{G((auto \textbf{AND} mixMotorRun \textbf{AND} X(auto \textbf{AND} rise(mixMotorComplete))) -> X(\textbf{NOT} mixMotorRun \textbf{AND} valveDRun))} \\
P19 & Mixer hold & The mixer keeps running until mixMotorComplete. & \texttt{G((auto \textbf{AND} mixMotorRun \textbf{AND} X(auto \textbf{AND} \textbf{NOT} mixMotorComplete)) -> X mixMotorRun)} \\
P20 & Sequencing & In automatic mode at most one material valve is open at a time, and the mixer never runs while a material valve is open. & \texttt{G(auto -> (\textbf{NOT}(valveARun \textbf{AND} valveBRun) \textbf{AND} \textbf{NOT}(valveARun \textbf{AND} valveCRun) \textbf{AND} \textbf{NOT}(valveBRun \textbf{AND} valveCRun) \textbf{AND} \textbf{NOT}(mixMotorRun \textbf{AND} fillOutputs)))} \\
P21 & Manual buttons ignored & In automatic mode the manual toggle buttons have no effect. & \texttt{G((auto \textbf{AND} X(auto \textbf{AND} \textbf{NOT} autoStart \textbf{AND} levelSensor < 50 \textbf{AND} \textbf{NOT} mixMotorComplete)) -> (valveDRun <-> X valveDRun))} (representative instance; analogous for each device/phase) \\
P22 & No restart & autoStart during a running sequence does not restart it. & \texttt{G((auto \textbf{AND} valveBRun \textbf{AND} X(auto \textbf{AND} rise(autoStart) \textbf{AND} levelSensor < 80)) -> X valveBRun)} \\
\bottomrule
\end{longtable}
}

\noindent\textbf{Assumptions and requirement gaps.}
{\small\begin{enumerate}
\item \textbf{Level thresholds} are read as "reaches" = \texttt{>=} (50, 80, 100 cm).
\item \textbf{End of the automatic process}: after valve D opens the requirement says "the process ends". We leave valve D open and do not constrain when it closes; a new autoStart is not required to close it.  Test cases therefore do not check anything after D opens except the single cycle in which it opens. \emph{(flagged as ambiguous)}
\item \textbf{Emergency stop release}: not specified whether the automatic sequence resumes. We assume it does \textbf{not} (P3), which is the safe interpretation; a new autoStart restarts from valve A.
\item \textbf{Mode switching}: when switching from automatic to manual, the automatic sequence is aborted; outputs are then under manual control. The reverse switch starts automatic mode idle. Not covered by test cases to avoid over-constraining.
\item \textbf{Manual toggle semantics} are explicitly rising-edge (requirement text), so holding a button for several cycles toggles once (P6).
\item \textbf{Mixer-complete signal while the mixer is not running} is ignored (follows from P9--P19).
\end{enumerate}}

\noindent\textbf{Validation.} reference implementation, simulation 6/6 cases PASS; mutation testing 2/2 mutants detected; 6 cases / 44 scan cycles.

\subsection{H02\_FlexProcess --- Flexible Processing Control}
\label{spec:H02_FlexProcess}

\noindent \texttt{H02\_FlexProcess}; 11 properties; source \texttt{specs/H02\_FlexProcess\_spec.md}.

{\small
\begin{longtable}{@{}p{0.05\linewidth}p{0.13\linewidth}p{0.30\linewidth}>{\raggedright\arraybackslash}p{0.42\linewidth}@{}}
\toprule
\textbf{ID} & \textbf{Category} & \textbf{Requirement} & \textbf{LTL} \\
\midrule
\endhead
P1 & Interlock & The conveyor never runs in both directions. & \texttt{G(\textbf{NOT} (conveyorLeft \textbf{AND} conveyorRight))} \\
P2 & Start & Motion starts only when raw material is at station 1 and the start button is pressed. & \texttt{G((\textbf{NOT} conveyorLeft \textbf{AND} \textbf{NOT} conveyorRight \textbf{AND} k = 0 \textbf{AND} X \textbf{NOT} (sensor1 \textbf{AND} station1Start)) -> X(\textbf{NOT} conveyorLeft \textbf{AND} \textbf{NOT} conveyorRight))} \\
P3 & Direction & The conveyor moves right when the target station number is greater than the current one. & \texttt{G((moving \textbf{AND} target > cur) -> conveyorRight)} \\
P4 & Direction & The conveyor moves left when the target station number is smaller than the current one. & \texttt{G((moving \textbf{AND} target < cur) -> conveyorLeft)} \\
P5 & Stop on arrival & The conveyor stops when the sensor of the target station detects the product. & \texttt{G((conveyorLeft \textbf{OR} conveyorRight) \textbf{AND} X sensor[target] -> X(\textbf{NOT} conveyorLeft \textbf{AND} \textbf{NOT} conveyorRight))} \\
P6 & Wait for completion & At a processing station the conveyor stays stopped until that station reports completion. & \texttt{G((at(target) \textbf{AND} \textbf{NOT} complete[target]) -> (\textbf{NOT} conveyorLeft \textbf{AND} \textbf{NOT} conveyorRight))} \\
P7 & Advance & On the completion signal of the current station the sequence advances to the next entry. & \texttt{G((at(target\_k) \textbf{AND} complete[target\_k]) -> X k = k + 1)} \\
P8 & Wrong signals & A completion signal from a station the product is not at has no effect. & \texttt{G((at(target) \textbf{AND} \textbf{NOT} complete[target]) -> X(\textbf{NOT} conveyorLeft \textbf{AND} \textbf{NOT} conveyorRight))} \\
P9 & Termination & After the final entry (station 5) is reached, the conveyor stops and the cycle ends. & \texttt{G(sensor5 \textbf{AND} k = processLength -> (\textbf{NOT} conveyorLeft \textbf{AND} \textbf{NOT} conveyorRight))} \\
P10 & Latched motion & While in transit the conveyor keeps running. & \texttt{G((conveyorRight \textbf{AND} X \textbf{NOT} sensor[target]) -> X conveyorRight)} and analogously for \texttt{conveyorLeft} \\
P11 & Sequence scope & Only entries 1..\texttt{processLength} are used. & the implementation never reads \texttt{processSequence[i]} for \texttt{i > processLength} \\
\bottomrule
\end{longtable}
}

\noindent\textbf{Assumptions and requirement gaps.}
{\small\begin{enumerate}
\item \textbf{The caller guarantees a valid sequence}: entries 2..4 for processing, the final entry is station 5, \texttt{processLength} 1..8 (stated in the requirement).
\item \textbf{Repeated station numbers} (e.g. [2,2,5]) are not specified; if the target equals the current station the product is already there and the completion signal advances the sequence immediately. Not tested.  \emph{(flagged as ambiguous)}
\item \textbf{Station 1 has no completion signal}; the start button plays that role.
\item \textbf{Sensors are FALSE between stations} (standard assumption of the line description).
\item \textbf{The start button is read as a level signal}; an edge-triggered implementation satisfies the same tests.
\end{enumerate}}

\noindent\textbf{Validation.} reference implementation, simulation 3/3 cases PASS; mutation testing 2/2 mutants detected; 3 cases / 21 scan cycles.

\subsection{H03\_RecipeMgr --- Recipe Management System}
\label{spec:H03_RecipeMgr}

\noindent \texttt{H03\_RecipeMgr}; 13 properties; source \texttt{specs/H03\_RecipeMgr\_spec.md}.

{\small
\begin{longtable}{@{}p{0.05\linewidth}p{0.13\linewidth}p{0.30\linewidth}>{\raggedright\arraybackslash}p{0.42\linewidth}@{}}
\toprule
\textbf{ID} & \textbf{Category} & \textbf{Requirement} & \textbf{LTL} \\
\midrule
\endhead
P1 & Query hit & Querying an existing number outputs that recipe without error. & \texttt{G((queryRecipe \textbf{AND} exists(recipeIn.recipeID)) -> (\textbf{NOT} error \textbf{AND} recipeQueryResult = db[slot(recipeIn.recipeID)]))} \\
P2 & Query miss & Querying a missing number reports 16\#8001. & \texttt{G((queryRecipe \textbf{AND} \textbf{NOT} exists(recipeIn.recipeID)) -> (error \textbf{AND} status = 16\#8001))} \\
P3 & Add duplicate & Adding an existing number reports 16\#8002 and does not modify the database. & \texttt{G((addRecipe \textbf{AND} exists(recipeIn.recipeID)) -> (error \textbf{AND} status = 16\#8002 \textbf{AND} X db = db))} \\
P4 & Add full & Adding with no free slot reports 16\#8003. & \texttt{G((addRecipe \textbf{AND} \textbf{NOT} exists(recipeIn.recipeID) \textbf{AND} noFreeSlot) -> (error \textbf{AND} status = 16\#8003 \textbf{AND} X db = db))} \\
P5 & Add success & A valid add stores the recipe in the free slot with the smallest index. & \texttt{G((addRecipe \textbf{AND} \textbf{NOT} exists(recipeIn.recipeID) \textbf{AND} hasFreeSlot) -> (recipeAdded \textbf{AND} X db[minFreeIndex] = recipeIn))} \\
P6 & Delete miss & Deleting a missing number reports 16\#8004. & \texttt{G((deleteRecipe \textbf{AND} \textbf{NOT} exists(recipeIn.recipeID)) -> (error \textbf{AND} status = 16\#8004))} \\
P7 & Delete success & Deleting an existing number clears that slot (\texttt{recipeID := 0}). & \texttt{G((deleteRecipe \textbf{AND} exists(recipeIn.recipeID)) -> (recipeDeleted \textbf{AND} X db[slot].recipeID = 0))} \\
P8 & Modify miss & Modifying a missing number reports 16\#8005. & \texttt{G((modifyRecipe \textbf{AND} \textbf{NOT} exists(recipeIn.recipeID)) -> (error \textbf{AND} status = 16\#8005))} \\
P9 & Modify success & Modifying an existing number replaces that slot with the input recipe. & \texttt{G((modifyRecipe \textbf{AND} exists(recipeIn.recipeID)) -> (recipeModified \textbf{AND} X db[slot] = recipeIn))} \\
P10 & Uniqueness & Recipe numbers in the database are never duplicated. & \texttt{G(FORALL i <> j: (db[i].recipeID <> 0) -> db[i].recipeID <> db[j].recipeID)} \\
P11 & No request & Without a request the database is unchanged. & \texttt{G((\textbf{NOT} addRecipe \textbf{AND} \textbf{NOT} deleteRecipe \textbf{AND} \textbf{NOT} modifyRecipe \textbf{AND} \textbf{NOT} queryRecipe) -> X db = db)} \\
P12 & Query read-only & A query never modifies the database. & \texttt{G(queryRecipe -> X db = db)} \\
P13 & Error status & \texttt{status = 16\#0000} whenever \texttt{error} is FALSE. & \texttt{G(\textbf{NOT} error -> status = 16\#0000)} \\
\bottomrule
\end{longtable}
}

\noindent\textbf{Assumptions and requirement gaps.}
{\small\begin{enumerate}
\item \textbf{Simultaneous requests} are not specified; the tests only ever assert one request at a time.  \emph{(flagged as ambiguous)}
\item \textbf{Request edge vs. level}: the requirement says "triggered"; the tests release the request between operations so that both readings behave identically.
\item \textbf{Success indicators} (\texttt{recipeAdded} etc.) are assumed to be TRUE in the cycle of a successful operation; their duration is not specified.
\item \textbf{\texttt{recipeQueryResult} after a failed query} is not specified and is not tested.
\item \textbf{Adding recipe number 0} is excluded, since 0 denotes an empty slot.
\end{enumerate}}

\noindent\textbf{Validation.} reference implementation, simulation 3/3 cases PASS; mutation testing 2/2 mutants detected; 3 cases / 23 scan cycles.

\subsection{H04\_FIFO --- FIFO First-In-First-Out Queue}
\label{spec:H04_FIFO}

\noindent \texttt{H04\_FIFO}; 12 properties; source \texttt{specs/H04\_FIFO\_spec.md}.

{\small
\begin{longtable}{@{}p{0.05\linewidth}p{0.13\linewidth}p{0.30\linewidth}>{\raggedright\arraybackslash}p{0.42\linewidth}@{}}
\toprule
\textbf{ID} & \textbf{Category} & \textbf{Requirement} & \textbf{LTL} \\
\midrule
\endhead
P1 & Enqueue & A rising edge of \texttt{enqueue} on a non-full queue appends \texttt{item} and increments the count. & \texttt{G((rise(enqueue) \textbf{AND} elementCount < cap \textbf{AND} \textbf{NOT} X reset) -> X elementCount = elementCount + 1)} \\
P2 & Dequeue & A rising edge of \texttt{dequeue} on a non-empty queue removes the front element, writes it to \texttt{item} and decrements the count. & \texttt{G((rise(dequeue) \textbf{AND} elementCount > 0 \textbf{AND} \textbf{NOT} X reset) -> (X elementCount = elementCount - 1 \textbf{AND} X item = frontValue))} \\
P3 & FIFO order & Values are removed in insertion order. & for any two enqueued values, the one enqueued first is dequeued first \\
P4 & Full & An enqueue into a full queue is rejected with 16\#8002 and changes nothing. & \texttt{G((rise(enqueue) \textbf{AND} elementCount = cap) -> (X status = 16\#8002 \textbf{AND} X elementCount = elementCount \textbf{AND} X item = item))} \\
P5 & Empty & A dequeue from an empty queue is rejected with 16\#8001 and changes nothing. & \texttt{G((rise(dequeue) \textbf{AND} elementCount = 0) -> (X status = 16\#8001 \textbf{AND} X elementCount = 0 \textbf{AND} X item = item))} \\
P6 & Reset & \texttt{reset} empties the queue and takes priority over both requests. & \texttt{G(reset -> (elementCount = 0 \textbf{AND} status = 16\#0000))} \\
P7 & Edge semantics & A held request performs only one operation. & \texttt{G((enqueue \textbf{AND} X enqueue) -> X elementCount = elementCount)} \\
P8 & Count bounds & The element count stays within 0..cap. & \texttt{G(elementCount >= 0 \textbf{AND} elementCount <= cap)} \\
P9 & Status & The status is 16\#0000 for idle cycles, resets and successful operations. & \texttt{G((\textbf{NOT} rejectedEnq \textbf{AND} \textbf{NOT} rejectedDeq) -> status = 16\#0000)} \\
P10 & Init & The queue is empty on the first invocation. & \texttt{elementCount = 0} (initial state) \\
P11 & Idle & Without a request the queue is unchanged. & \texttt{G((\textbf{NOT} rise(enqueue) \textbf{AND} \textbf{NOT} rise(dequeue) \textbf{AND} \textbf{NOT} reset) -> X elementCount = elementCount)} \\
P12 & Wrap-around & FIFO order is preserved across the physical end of the buffer. & contained in P3 \\
\bottomrule
\end{longtable}
}

\noindent\textbf{Assumptions and requirement gaps.}
{\small\begin{enumerate}
\item \textbf{The caller never requests enqueue and dequeue in the same invocation} (stated in the requirement).
\item \textbf{Reset priority} over both operations is explicit; buffer contents need not be cleared.
\item \textbf{\texttt{item} on a rejected operation} stays unchanged (P4, P5).
\item \textbf{Capacity in the tests is 8} (\texttt{ARRAY[0..7]}), matching the project's reference implementation; the task itself derives capacity from the caller-supplied buffer, which the harness cannot represent as \texttt{VARIANT}.  \emph{(flagged as ambiguous)}
\item \textbf{Status of an idle cycle} is 16\#0000 (P9).
\end{enumerate}}

\noindent\textbf{Validation.} reference implementation, simulation 5/5 cases PASS; mutation testing 2/2 mutants detected; 5 cases / 66 scan cycles.

\subsection{H05\_MechArm --- Mechanical Arm Control}
\label{spec:H05_MechArm}

\noindent \texttt{H05\_MechArm}; 19 properties; source \texttt{specs/H05\_MechArm\_spec.md}.

{\small
\begin{longtable}{@{}p{0.05\linewidth}p{0.13\linewidth}p{0.30\linewidth}>{\raggedright\arraybackslash}p{0.42\linewidth}@{}}
\toprule
\textbf{ID} & \textbf{Category} & \textbf{Requirement} & \textbf{LTL} \\
\midrule
\endhead
P1 & Emergency stop & While the emergency stop is active every output is OFF. & \texttt{G(estop -> allOff)} \\
P2 & Limit switches & No motion/actuation is commanded into a position that is already reached. & \texttt{G((DownLimitSwitch -> \textbf{NOT} MoveDown) \textbf{AND} (UpLimitSwitch -> \textbf{NOT} MoveUp) \textbf{AND} (LeftLimitSwitch -> \textbf{NOT} MoveLeft) \textbf{AND} (RightLimitSwitch -> \textbf{NOT} MoveRight) \textbf{AND} (GripLimitSwitch -> \textbf{NOT} ActivateGrip) \textbf{AND} (ReleaseLimitSwitch -> \textbf{NOT} ActivateRelease))} \\
P3 & Opposite outputs & Opposite commands are never active together. & \texttt{G(\textbf{NOT}(MoveDown \textbf{AND} MoveUp) \textbf{AND} \textbf{NOT}(MoveLeft \textbf{AND} MoveRight) \textbf{AND} \textbf{NOT}(ActivateGrip \textbf{AND} ActivateRelease))} \\
P4 & Init & All outputs are OFF at power-up. & \texttt{allOff} (initial state) \\
P5 & No resume after estop & After the emergency stop is released, automatic motion does not resume by itself. & \texttt{G((estop \textbf{AND} X(auto \textbf{AND} \textbf{NOT} Start)) -> X allOff)} \\
P6 & Jog & Each output equals its button AND NOT its limit switch (shown for down; analogous for up/left/right/grip/release). & \texttt{G((EmergencyStop \textbf{AND} ManualMode) -> (MoveDown <-> (DownButton \textbf{AND} \textbf{NOT} DownLimitSwitch)))} \\
P7 & Start ignored & Start has no effect in manual mode. & implied by P6 (outputs depend only on buttons and switches) \\
P8 & Wait for Start & When idle, nothing moves until Start is pressed. & \texttt{G((auto \textbf{AND} allOff \textbf{AND} X(auto \textbf{AND} \textbf{NOT} Start)) -> X allOff)} \\
P9 & First action & From home, Start begins the down movement. & \texttt{G((auto \textbf{AND} allOff \textbf{AND} home \textbf{AND} X(auto \textbf{AND} rise(Start))) -> X MoveDown)} \\
P10 & Manual buttons ignored & Jog buttons have no effect in automatic modes. & \texttt{G((auto \textbf{AND} allOff \textbf{AND} X(auto \textbf{AND} \textbf{NOT} Start \textbf{AND} (DownButton \textbf{OR} UpButton \textbf{OR} LeftButton \textbf{OR} RightButton \textbf{OR} GripButton \textbf{OR} ReleaseButton))) -> X allOff)} \\
P11 & Action completes at limit & A running action stops when its limit switch is reached (down shown). & \texttt{G((auto \textbf{AND} MoveDown \textbf{AND} X(auto \textbf{AND} DownLimitSwitch)) -> X \textbf{NOT} MoveDown)} \\
P12 & Action persists & A running action continues until its limit switch is reached (down shown). & \texttt{G((auto \textbf{AND} MoveDown \textbf{AND} X(auto \textbf{AND} \textbf{NOT} DownLimitSwitch)) -> X MoveDown)} \\
P13 & Order & Actions follow the fixed order: grip is only commanded after down was completed, up after grip, right after up, \ldots{} (representative instance). & \texttt{G(ActivateGrip -> DownLimitSwitch)}, \texttt{G(MoveRight -> (UpLimitSwitch \textbf{AND} GripLimitSwitch))}, \texttt{G(ActivateRelease -> (DownLimitSwitch \textbf{AND} RightLimitSwitch))}, \texttt{G(MoveLeft -> (UpLimitSwitch \textbf{AND} ReleaseLimitSwitch))} \\
P14 & One action per Start & After an action completes, the arm stays stopped until the next Start. & \texttt{G((SingleStepMode \textbf{AND} EmergencyStop \textbf{AND} allOff \textbf{AND} X(SingleStepMode \textbf{AND} EmergencyStop \textbf{AND} \textbf{NOT} Start)) -> X allOff)} \\
P15 & Start advances & When stopped after step k, Start begins step k+1 (e.g. at the bottom-left gripped position, Start $\rightarrow$ MoveUp). & \texttt{G((SingleStepMode \textbf{AND} EmergencyStop \textbf{AND} allOff \textbf{AND} DownLimitSwitch \textbf{AND} GripLimitSwitch \textbf{AND} LeftLimitSwitch \textbf{AND} X(SingleStepMode \textbf{AND} Start)) -> X MoveUp)} (one instance per step) \\
P16 & Held Start & Holding Start does not start the following action (edge semantics). & \texttt{G((SingleStepMode \textbf{AND} Start \textbf{AND} X(SingleStepMode \textbf{AND} Start) \textbf{AND} X allOff\_after\_completion) -> X allOff)} \\
P17 & Chaining & When an action completes, the next action starts without a new Start (down$\rightarrow$grip shown). & \texttt{G((SingleCycleMode \textbf{AND} EmergencyStop \textbf{AND} MoveDown \textbf{AND} X(SingleCycleMode \textbf{AND} DownLimitSwitch \textbf{AND} \textbf{NOT} GripLimitSwitch)) -> X F ActivateGrip)} (within the next cycle) \\
P18 & Stop at home & After the left movement completes (arm back at home) the arm stops and waits for Start. & \texttt{G((SingleCycleMode \textbf{AND} MoveLeft \textbf{AND} X(SingleCycleMode \textbf{AND} LeftLimitSwitch)) -> X allOff)} and P8 afterwards \\
P19 & Repeat & After the left movement completes, the next cycle's down movement starts without Start. & \texttt{G((ContinuousMode \textbf{AND} MoveLeft \textbf{AND} X(ContinuousMode \textbf{AND} LeftLimitSwitch \textbf{AND} \textbf{NOT} DownLimitSwitch)) -> X F MoveDown)} (within the next cycle) \\
\bottomrule
\end{longtable}
}

\noindent\textbf{Assumptions and requirement gaps.}
{\small\begin{enumerate}
\item \textbf{Start is edge-triggered} in single-step mode ("pressed again"); holding Start must not run several steps. In single-cycle/continuous modes Start only matters when idle.
\item \textbf{Transition timing}: when a limit switch is reached, the completed action is OFF in that cycle; the next action may start in the same cycle or in the following one. Test cases allow one cycle of slack (two consecutive steps with the same inputs).  \emph{(flagged as ambiguous)}
\item \textbf{Emergency stop release}: not specified whether the interrupted cycle resumes. We assume it does \textbf{not} (P5); the sequence restarts from the beginning with a new Start (the arm is assumed to be returned manually to home).
\item \textbf{Mode change while running}: not specified; not tested.
\item \textbf{Limit switches} are assumed to be exclusive per axis (up/down, left/right, grip/release) and TRUE only at the end positions.
\item \textbf{Initial position} is top-left with the gripper released, as stated.
\end{enumerate}}

\noindent\textbf{Validation.} reference implementation, simulation 7/7 cases PASS; mutation testing 2/2 mutants detected; 7 cases / 135 scan cycles.

\subsection{H06\_Production --- Production Control}
\label{spec:H06_Production}

\noindent \texttt{H06\_Production}; 21 properties; source \texttt{specs/H06\_Production\_spec.md}.

{\small
\begin{longtable}{@{}p{0.05\linewidth}p{0.13\linewidth}p{0.30\linewidth}>{\raggedright\arraybackslash}p{0.42\linewidth}@{}}
\toprule
\textbf{ID} & \textbf{Category} & \textbf{Requirement} & \textbf{LTL} \\
\midrule
\endhead
P1 & Interlock & Forward and reverse are never commanded together. & \texttt{G(\textbf{NOT} (MotorForward \textbf{AND} MotorReverse))} \\
P2 & Init & All Boolean outputs OFF and firingTemp = 0 at power-up. & \texttt{stopped \textbf{AND} \textbf{NOT} kick \textbf{AND} \textbf{NOT} CompletionLight \textbf{AND} firingTemp = 0.0} (initial state) \\
P3 & Jog & Forward/reverse buttons jog the motor with interlock; no button $\rightarrow$ stop. & \texttt{G(\textbf{NOT} Mode -> ((MotorForward <-> (ForwardButton \textbf{AND} \textbf{NOT} ReverseButton)) \textbf{AND} (MotorReverse <-> (ReverseButton \textbf{AND} \textbf{NOT} ForwardButton))))} \\
P4 & No process outputs  & kick and CompletionLight are OFF in manual mode. & \texttt{G(\textbf{NOT} Mode -> (\textbf{NOT} kick \textbf{AND} \textbf{NOT} CompletionLight))} \\
P5 & No reverse & The motor never reverses in automatic mode. & \texttt{G(Mode -> \textbf{NOT} MotorReverse)} \\
P6 & Start & Material at the feed station (line idle) starts forward motion. & \texttt{G((Mode \textbf{AND} stopped \textbf{AND} \textbf{NOT} CompletionLight \textbf{AND} \textbf{NOT} anySensor \textbf{AND} X(Mode \textbf{AND} SensorFeed)) -> X MotorForward)} \\
P7 & Stop on arrival & Reaching any station sensor while moving stops the motor. & \texttt{G((Mode \textbf{AND} MotorForward \textbf{AND} X(Mode \textbf{AND} anySensor)) -> X stopped)} \\
P8 & Latched motion & Between stations the motor keeps running. & \texttt{G((Mode \textbf{AND} MotorForward \textbf{AND} X(Mode \textbf{AND} \textbf{NOT} anySensor)) -> X MotorForward)} \\
P9 & Wash $\rightarrow$ weigh & At the wash station, washComplete restarts forward motion. & \texttt{G((Mode \textbf{AND} SensorWash \textbf{AND} stopped \textbf{AND} X(Mode \textbf{AND} SensorWash \textbf{AND} rise(washComplete))) -> X MotorForward)} \\
P10 & Wait at wash & Without washComplete the motor stays stopped at the wash station. & \texttt{G((Mode \textbf{AND} SensorWash \textbf{AND} stopped \textbf{AND} X(Mode \textbf{AND} SensorWash \textbf{AND} \textbf{NOT} washComplete)) -> X stopped)} \\
P11 & Qualified & At the weigh station, weightComplete with a qualified weight: no kick, motor forward, firingTemp = 500 + (weight $-$ 100) $\times$ 0.5. & \texttt{G((Mode \textbf{AND} SensorWeight \textbf{AND} stopped \textbf{AND} X(Mode \textbf{AND} SensorWeight \textbf{AND} rise(weightComplete) \textbf{AND} qualified)) -> X(\textbf{NOT} kick \textbf{AND} MotorForward \textbf{AND} firingTemp = 500.0 + (weight - 100.0) * 0.5))} \\
P12 & Rejected & At the weigh station, weightComplete with an unqualified weight: kick is issued and the motor does not move forward. & \texttt{G((Mode \textbf{AND} SensorWeight \textbf{AND} stopped \textbf{AND} X(Mode \textbf{AND} SensorWeight \textbf{AND} rise(weightComplete) \textbf{AND} \textbf{NOT} qualified)) -> X(kick \textbf{AND} \textbf{NOT} MotorForward))} \\
P13 & Kick only at weigh & kick is never issued except in the cycle of a weightComplete at the weigh station with an unqualified weight. & \texttt{G(kick -> (Mode \textbf{AND} SensorWeight \textbf{AND} weightComplete \textbf{AND} \textbf{NOT} qualified))} \\
P14 & Temperature range & firingTemp is always within 500..550 once set, or 0 before any qualified weighing. & \texttt{G(firingTemp = 0.0 \textbf{OR} (firingTemp >= 500.0 \textbf{AND} firingTemp <= 550.0))} \\
P15 & Temperature hold & firingTemp changes only at a qualified weighing. & \texttt{G((Mode \textbf{AND} \textbf{NOT} (SensorWeight \textbf{AND} X rise(weightComplete))) -> (firingTemp = X firingTemp))} \\
P16 & Firing $\rightarrow$ output & At the firing station, firingComplete restarts forward motion. & \texttt{G((Mode \textbf{AND} SensorFiring \textbf{AND} stopped \textbf{AND} X(Mode \textbf{AND} SensorFiring \textbf{AND} firingComplete)) -> X MotorForward)} \\
P17 & Light ON & Arrival at the output station turns the completion light on. & \texttt{G((Mode \textbf{AND} MotorForward \textbf{AND} X(Mode \textbf{AND} SensorOutput)) -> X(CompletionLight \textbf{AND} stopped))} \\
P18 & Light hold & The light stays on while the product remains at the output station. & \texttt{G((Mode \textbf{AND} CompletionLight \textbf{AND} X(Mode \textbf{AND} SensorOutput)) -> X CompletionLight)} \\
P19 & Light OFF & When the product is removed the light goes off. & \texttt{G((Mode \textbf{AND} X(Mode \textbf{AND} \textbf{NOT} SensorOutput)) -> X \textbf{NOT} CompletionLight)} \\
P20 & Wrong-station signals & A complete signal at a station other than its own has no effect on the motor. & e.g. \texttt{G((Mode \textbf{AND} SensorWash \textbf{AND} stopped \textbf{AND} X(Mode \textbf{AND} SensorWash \textbf{AND} \textbf{NOT} washComplete \textbf{AND} (weightComplete \textbf{OR} firingComplete))) -> X stopped)} and analogously for the weigh and firing stations \\
P21 & Manual buttons ignored & Jog buttons have no effect in automatic mode. & \texttt{G((Mode \textbf{AND} stopped \textbf{AND} \textbf{NOT} SensorFeed \textbf{AND} X(Mode \textbf{AND} ForwardButton \textbf{AND} \textbf{NOT} SensorFeed \textbf{AND} \textbf{NOT} washComplete \textbf{AND} \textbf{NOT} weightComplete \textbf{AND} \textbf{NOT} firingComplete)) -> X stopped)} \\
\bottomrule
\end{longtable}
}

\noindent\textbf{Assumptions and requirement gaps.}
{\small\begin{enumerate}
\item \textbf{P4}: the requirement does not mention kick/light in manual mode; we assume OFF.
\item \textbf{kick duration}: not specified. We treat kick as a one-cycle pulse issued in the cycle the weight is judged unqualified (P12/P13); the tests only check kick = 1 in that cycle and kick = 0 after the rejected material has left the weigh station. An implementation that holds kick while SensorWeight remains TRUE also satisfies the tests except step 5 of TC04 (material already removed).  \emph{(flagged as ambiguous)}
\item \textbf{After rejection}: the requirement does not say what happens next. We assume the rejected material is removed by the kicker and the line returns to idle, ready for the next material at the feed station (TC04 step 6).
\item \textbf{Sensor FALSE between stations}: assumed (sensors only see the material at the station).
\item \textbf{Mode switch during motion}: switching to manual stops automatic motion immediately; switching back starts automatic mode idle (TC06). Not constrained beyond that.
\item \textbf{Temperature formula} \texttt{500 + (weight $-$ 100) $\times$ 0.5} follows from the linear map 100..200 $\rightarrow$ 500..550; at weight = 100 $\rightarrow$ 500.0, 150 $\rightarrow$ 525.0, 200 $\rightarrow$ 550.0. REAL outputs are compared exactly in the tests because these values are exactly representable.
\end{enumerate}}

\noindent\textbf{Validation.} reference implementation, simulation 6/6 cases PASS; mutation testing 2/2 mutants detected; 6 cases / 95 scan cycles.

\subsection{M01\_PumpSwitch --- Multi-Pump Switching Control System}
\label{spec:M01_PumpSwitch}

\noindent \texttt{M01\_PumpSwitch}; 8 properties; source \texttt{specs/M01\_PumpSwitch\_spec.md}.

{\small
\begin{longtable}{@{}p{0.05\linewidth}p{0.13\linewidth}p{0.30\linewidth}>{\raggedright\arraybackslash}p{0.42\linewidth}@{}}
\toprule
\textbf{ID} & \textbf{Category} & \textbf{Requirement} & \textbf{LTL} \\
\midrule
\endhead
P1 & Stop & \texttt{stop} stops every pump in the same cycle. & \texttt{G(stop -> FORALL i: \textbf{NOT} pumpGroup[i].RunComd)} \\
P2 & Idle & Before any \texttt{start}, no pump runs. & \texttt{G(\textbf{NOT} running -> FORALL i: \textbf{NOT} pumpGroup[i].RunComd)} \\
P3 & Manual & In manual mode a started pump runs exactly if it is selected. & \texttt{G((running \textbf{AND} \textbf{NOT} mode) -> FORALL i: (pumpGroup[i].RunComd <-> pumpGroup[i].selection))} \\
P4 & Automatic & In automatic mode exactly three pumps run. & \texttt{G((running \textbf{AND} mode) -> count(pumpGroup[i].RunComd) = 3)} \\
P5 & Automatic & A running pump has a priority at least as high as every stopped pump. & \texttt{G((running \textbf{AND} mode \textbf{AND} pumpGroup[i].RunComd \textbf{AND} \textbf{NOT} pumpGroup[j].RunComd) -> pumpGroup[i].priority >= pumpGroup[j].priority)} \\
P6 & Automatic & The manual selection has no effect in automatic mode. & follows from P4/P5 (outputs depend only on priority) \\
P7 & Read-only & The function block never writes \texttt{priority} or \texttt{selection}. & \texttt{G(X pumpGroup[i].priority = pumpGroup[i].priority \textbf{AND} X pumpGroup[i].selection = pumpGroup[i].selection)} \\
P8 & Standby & Exactly two pumps are on standby in automatic mode. & equivalent to P4 for a five-pump group \\
\bottomrule
\end{longtable}
}

\noindent\textbf{Assumptions and requirement gaps.}
{\small\begin{enumerate}
\item \textbf{Start/stop semantics}  The requirement describes buttons without fixing latching. The team adopted a \textbf{level} reading: \texttt{RunComd := start AND \ldots{}}, so pumps run only while \texttt{start} is held and \texttt{stop} clears everything. \emph{(flagged as ambiguous)}
\item \textbf{Priority ties}: not specified. Any three pumps with the highest priorities satisfy P4/P5; with all priorities equal, exactly three must run (TC03).
\item \textbf{Selection changes while running} take effect immediately in manual mode (TC01).
\item \textbf{Mode switching while running} re-evaluates the outputs according to the new mode (TC03).
\end{enumerate}}

\noindent\textbf{Validation.} reference implementation, simulation 3/3 cases PASS; mutation testing 2/2 mutants detected; 3 cases / 15 scan cycles.

\subsection{M02\_AnalogBatch --- Analog Batch Processing}
\label{spec:M02_AnalogBatch}

\noindent \texttt{M02\_AnalogBatch}; 12 properties; source \texttt{specs/M02\_AnalogBatch\_spec.md}.

{\small
\begin{longtable}{@{}p{0.05\linewidth}p{0.13\linewidth}p{0.30\linewidth}>{\raggedright\arraybackslash}p{0.42\linewidth}@{}}
\toprule
\textbf{ID} & \textbf{Category} & \textbf{Requirement} & \textbf{LTL} \\
\midrule
\endhead
P1 & Scope & Only the first \texttt{count} channels are processed; the rest are untouched. & \texttt{G(i > count -> (X ch[i].outputValue = ch[i].outputValue \textbf{AND} X ch[i].retVal = ch[i].retVal))} \\
P2 & Enable & A disabled channel is not converted. & \texttt{G(\textbf{NOT} ch[i].enable -> X ch[i].outputValue = ch[i].outputValue)} \\
P3 & Mode error & An invalid \texttt{measuringMode} yields 16\#8001 and no conversion. & \texttt{G((ch[i].enable AND ch[i].measuringMode NOT IN \{0,1,2\}) -> (ch[i].retVal = 16\#8001 AND X ch[i].outputValue = ch[i].outputValue))} \\
P4 & Limit error & \texttt{hiLIM <= loLIM} yields 16\#8002 and no conversion. & \texttt{G((ch[i].enable \textbf{AND} ch[i].hiLIM <= ch[i].loLIM) -> (ch[i].retVal = 16\#8002 \textbf{AND} X ch[i].outputValue = ch[i].outputValue))} \\
P5 & Unipolar & Mode 0 with \texttt{bipolar = FALSE} maps 0..27648 linearly to loLIM..hiLIM. & \texttt{G(valid0(i) \textbf{AND} \textbf{NOT} ch[i].bipolar -> ch[i].outputValue = ch[i].loLIM + (analogValues[i] / 27648) * (ch[i].hiLIM - ch[i].loLIM))} \\
P6 & Bipolar & Mode 0 with \texttt{bipolar = TRUE} maps $-$27648..27648 linearly to loLIM..hiLIM. & \texttt{G(valid0(i) \textbf{AND} ch[i].bipolar -> ch[i].outputValue = ch[i].loLIM + ((analogValues[i] + 27648) / 55296) * (ch[i].hiLIM - ch[i].loLIM))} \\
P7 & Standard temp & Mode 1 gives one tenth of the raw value. & \texttt{G(valid1(i) -> ch[i].outputValue = analogValues[i] / 10.0)} \\
P8 & Climatic temp & Mode 2 gives one hundredth of the raw value. & \texttt{G(valid2(i) -> ch[i].outputValue = analogValues[i] / 100.0)} \\
P9 & Success status & A successfully converted channel reports 16\#0000. & \texttt{G(converted(i) -> ch[i].retVal = 16\#0000)} \\
P10 & Read-only input & \texttt{analogValues} is never modified. & \texttt{G(X analogValues[i] = analogValues[i])} \\
P11 & Read-only config & Configuration fields are never modified. & \texttt{G(X ch[i].enable = ch[i].enable \textbf{AND} X ch[i].hiLIM = ch[i].hiLIM \textbf{AND} X ch[i].loLIM = ch[i].loLIM \textbf{AND} X ch[i].bipolar = ch[i].bipolar \textbf{AND} X ch[i].measuringMode = ch[i].measuringMode)} \\
P12 & Independence & An error on one channel does not prevent other channels from being converted. & \texttt{G((ch[i].retVal = 16\#8001 \textbf{OR} ch[i].retVal = 16\#8002) \textbf{AND} valid0(j) -> converted(j))} \\
\bottomrule
\end{longtable}
}

\noindent\textbf{Assumptions and requirement gaps.}
{\small\begin{enumerate}
\item \textbf{Limit check scope}: the requirement states the \texttt{hiLIM <= loLIM} error without restricting it to mode 0; we apply it to every enabled channel.  \emph{(flagged as ambiguous)}
\item \textbf{Error priority}: an invalid mode (16\#8001) is checked before the limits (16\#8002); the requirement lists them in that order.
\item \textbf{\texttt{retVal} of disabled channels} is not specified and is not tested.
\item \textbf{Out-of-range raw values} (beyond $\pm$27648) are not specified; linear extrapolation is assumed and not tested.
\end{enumerate}}

\noindent\textbf{Validation.} reference implementation, simulation 3/3 cases PASS; mutation testing 2/2 mutants detected; 3 cases / 11 scan cycles.

\subsection{M03\_MovingAvg --- Calculate Moving Average}
\label{spec:M03_MovingAvg}

\noindent \texttt{M03\_MovingAvg}; 11 properties; source \texttt{specs/M03\_MovingAvg\_spec.md}.

{\small
\begin{longtable}{@{}p{0.05\linewidth}p{0.13\linewidth}p{0.30\linewidth}>{\raggedright\arraybackslash}p{0.42\linewidth}@{}}
\toprule
\textbf{ID} & \textbf{Category} & \textbf{Requirement} & \textbf{LTL} \\
\midrule
\endhead
P1 & Window error & A window size below 1 or above 100 produces an error and 16\#8200. & \texttt{G((windowSize < 1 \textbf{OR} windowSize > 100) -> (error \textbf{AND} status = 16\#8200))} \\
P2 & Normal status & With a valid window size no error is reported. & \texttt{G((windowSize >= 1 \textbf{AND} windowSize <= 100 \textbf{AND} \textbf{NOT} error\_other) -> (\textbf{NOT} error \textbf{AND} status = 16\#0000))} \\
P3 & Cyclic read & With \texttt{cyclicExecution} TRUE a new value is read every cycle. & \texttt{G((cyclicExecution \textbf{AND} validWindow \textbf{AND} \textbf{NOT} reset) -> X n = min(n + 1, windowSize))} \\
P4 & Trigger read & With \texttt{cyclicExecution} FALSE a value is read exactly on each rising edge of \texttt{trigger}. & \texttt{G((\textbf{NOT} cyclicExecution \textbf{AND} \textbf{NOT} trigger \textbf{AND} X trigger \textbf{AND} validWindow \textbf{AND} \textbf{NOT} reset) -> X n = min(n + 1, windowSize))} \\
P5 & No read & Without a cyclic read or a trigger edge the stored data is unchanged. & \texttt{G((\textbf{NOT} cyclicExecution \textbf{AND} \textbf{NOT} (\textbf{NOT} trigger \textbf{AND} X trigger)) -> X average = average)} \\
P6 & Trigger ignored & \texttt{cyclicExecution} overrides \texttt{trigger}. & implied by P3 \\
P7 & Average & The output equals the arithmetic mean of the stored window. & \texttt{G((validWindow \textbf{AND} n > 0) -> average = sum(window) / n)} \\
P8 & Window reached & \texttt{windowSizeReached} is TRUE exactly when the window is full. & \texttt{G(windowSizeReached <-> (n >= windowSize))} \\
P9 & Reset & \texttt{reset} clears the calculation. & \texttt{G(reset -> (average = 0 \textbf{AND} \textbf{NOT} windowSizeReached))} \\
P10 & Replacement & Once full, each new value replaces the oldest one. & \texttt{G((windowSizeReached \textbf{AND} read) -> X n = windowSize)} \\
P11 & Init & At power-up the average is 0 and the window is not reached. & \texttt{average = 0 \textbf{AND} \textbf{NOT} windowSizeReached} (initial state) \\
\bottomrule
\end{longtable}
}

\noindent\textbf{Assumptions and requirement gaps.}
{\small\begin{enumerate}
\item \textbf{Average before the window is full}: the requirement does not specify it; we assume the mean of the samples collected so far (P7), which is the usual behaviour.  \emph{(flagged as ambiguous)}
\item \textbf{Reset priority}: \texttt{reset} is evaluated before reading a new value; the window-size error is checked before \texttt{reset}.
\item \textbf{Status on window error} is 16\#8200 and the average is set to 0.
\item \textbf{Comparison tolerance}: \texttt{average} is LREAL; tests use values whose means are exactly representable except TC01 step 3 (36.666\ldots{}), which the runner compares with full double precision.
\end{enumerate}}

\noindent\textbf{Validation.} reference implementation, simulation 5/5 cases PASS; mutation testing 2/2 mutants detected; 5 cases / 24 scan cycles.

\subsection{M04\_TempCtrl --- Temperature Logic Control}
\label{spec:M04_TempCtrl}

\noindent \texttt{M04\_TempCtrl}; 12 properties; source \texttt{specs/M04\_TempCtrl\_spec.md}.

{\small
\begin{longtable}{@{}p{0.05\linewidth}p{0.13\linewidth}p{0.30\linewidth}>{\raggedright\arraybackslash}p{0.42\linewidth}@{}}
\toprule
\textbf{ID} & \textbf{Category} & \textbf{Requirement} & \textbf{LTL} \\
\midrule
\endhead
P1 & Display & The displays mirror the inputs every cycle. & \texttt{G(CurrentTemperature = TemperatureSensorInput \textbf{AND} SetTempDisplay = SetTemperature)} \\
P2 & Status & \texttt{HeaterStatus} reflects the heater output. & \texttt{G(HeaterStatus <-> HeaterOutput)} \\
P3 & Sensor fault & A value below 0 or above 100 turns the heater off. & \texttt{G((TemperatureSensorInput < 0.0 \textbf{OR} TemperatureSensorInput > 100.0) -> \textbf{NOT} HeaterOutput)} \\
P4 & Heating & Below the setpoint, without a fault or protection, the heater is on. & \texttt{G((TemperatureSensorInput < SetTemperature \textbf{AND} \textbf{NOT} fault \textbf{AND} \textbf{NOT} OverheatProtection) -> HeaterOutput)} \\
P5 & Cooling & At or above the setpoint the heater is off. & \texttt{G(TemperatureSensorInput >= SetTemperature -> \textbf{NOT} HeaterOutput)} \\
P6 & Overheat trip & After 5 s of continuous heating, protection is activated. & \texttt{G(heatAcc >= 5s -> OverheatProtection)} \\
P7 & Protection action & While protection is active the heater is off. & \texttt{G(OverheatProtection -> \textbf{NOT} HeaterOutput)} \\
P8 & Protection hold & Protection is held for at least 3 s. & \texttt{G(OverheatProtection -> (protAcc >= 3s \textbf{OR} X OverheatProtection))} \\
P9 & Protection release & Protection is released once 3 s have elapsed. & \texttt{G(protAcc >= 3s -> \textbf{NOT} X OverheatProtection)} \\
P10 & No false trip & Protection is not activated before 5 s of continuous heating. & \texttt{G((heatAcc < 5s \textbf{AND} \textbf{NOT} OverheatProtection) -> \textbf{NOT} X OverheatProtection)} \\
P11 & Timer restart & Interrupting the heater restarts the 5 s measurement. & \texttt{G(\textbf{NOT} HeaterOutput -> X heatAcc <= CycleTime)} \\
P12 & Init & At power-up the heater and protection are off. & \texttt{\textbf{NOT} HeaterOutput \textbf{AND} \textbf{NOT} OverheatProtection} (initial state) \\
\bottomrule
\end{longtable}
}

\noindent\textbf{Assumptions and requirement gaps.}
{\small\begin{enumerate}
\item \textbf{Equality with the setpoint}: "lower than" and "higher than" leave \texttt{=} undefined; we treat \texttt{temperature = setpoint} as "do not heat" (P5).  \emph{(flagged as ambiguous)}
\item \textbf{Valid sensor range} is inclusive: 0 and 100 are valid values (the fault condition is \texttt{< 0} or \texttt{> 100}).
\item \textbf{Protection during a sensor fault}: the heater is off in both cases; no interaction is specified.
\item \textbf{After protection is released} normal control resumes immediately if the temperature is still below the setpoint.
\item The displays are updated even during a sensor fault (P1).
\end{enumerate}}

\noindent\textbf{Validation.} reference implementation, simulation 4/4 cases PASS; mutation testing 2/2 mutants detected; 4 cases / 22 scan cycles.

\subsection{M05\_ProdLine --- Automated Production Line Control}
\label{spec:M05_ProdLine}

\noindent \texttt{M05\_ProdLine}; 18 properties; source \texttt{specs/M05\_ProdLine\_spec.md}.

{\small
\begin{longtable}{@{}p{0.05\linewidth}p{0.13\linewidth}p{0.30\linewidth}>{\raggedright\arraybackslash}p{0.42\linewidth}@{}}
\toprule
\textbf{ID} & \textbf{Category} & \textbf{Requirement} & \textbf{LTL} \\
\midrule
\endhead
P1 & Interlock & Forward and reverse are never commanded together. & \texttt{G(\textbf{NOT} (MotorForward \textbf{AND} MotorReverse))} \\
P2 & Init & All outputs are OFF at power-up. & \texttt{\textbf{NOT} MotorForward \textbf{AND} \textbf{NOT} MotorReverse \textbf{AND} \textbf{NOT} CompletionLight} (initial state) \\
P3 & Jog forward & Forward button alone drives the motor forward. & \texttt{G((\textbf{NOT} Mode \textbf{AND} ForwardButton \textbf{AND} \textbf{NOT} ReverseButton) -> MotorForward)} \\
P4 & Jog reverse & Reverse button alone drives the motor in reverse. & \texttt{G((\textbf{NOT} Mode \textbf{AND} ReverseButton \textbf{AND} \textbf{NOT} ForwardButton) -> MotorReverse)} \\
P5 & Interlock & Both buttons pressed: neither direction runs. & \texttt{G((\textbf{NOT} Mode \textbf{AND} ForwardButton \textbf{AND} ReverseButton) -> (\textbf{NOT} MotorForward \textbf{AND} \textbf{NOT} MotorReverse))} \\
P6 & Stop on release & No button: motor stops (no latching in manual). & \texttt{G((\textbf{NOT} Mode \textbf{AND} \textbf{NOT} ForwardButton \textbf{AND} \textbf{NOT} ReverseButton) -> (\textbf{NOT} MotorForward \textbf{AND} \textbf{NOT} MotorReverse))} \\
P7 & Light & The completion light is not used in manual mode.  & \texttt{G(\textbf{NOT} Mode -> \textbf{NOT} CompletionLight)} \\
P8 & No reverse & The motor never reverses in automatic mode. & \texttt{G(Mode -> \textbf{NOT} MotorReverse)} \\
P9 & Start at A & Product at A and ButtonA pressed (while idle) start forward motion. & \texttt{G((Mode \textbf{AND} \textbf{NOT} MotorForward \textbf{AND} \textbf{NOT} CompletionLight \textbf{AND} X(Mode \textbf{AND} SensorA \textbf{AND} rise(ButtonA))) -> X MotorForward)} \\
P10 & Stop at B & Arrival at B stops the motor. & \texttt{G((Mode \textbf{AND} MotorForward \textbf{AND} X(Mode \textbf{AND} SensorB)) -> X \textbf{NOT} MotorForward)} \\
P11 & Continue from B & At B, ButtonB restarts forward motion. & \texttt{G((Mode \textbf{AND} SensorB \textbf{AND} \textbf{NOT} MotorForward \textbf{AND} X(Mode \textbf{AND} SensorB \textbf{AND} rise(ButtonB))) -> X MotorForward)} \\
P12 & Stop at C & Arrival at C stops the motor. & \texttt{G((Mode \textbf{AND} MotorForward \textbf{AND} X(Mode \textbf{AND} SensorC)) -> X \textbf{NOT} MotorForward)} \\
P13 & Light ON & At C, ButtonC turns the completion light on. & \texttt{G((Mode \textbf{AND} SensorC \textbf{AND} \textbf{NOT} MotorForward \textbf{AND} X(Mode \textbf{AND} SensorC \textbf{AND} rise(ButtonC))) -> X CompletionLight)} \\
P14 & Light OFF & When the product is removed from C the light goes off. & \texttt{G((Mode \textbf{AND} X(Mode \textbf{AND} \textbf{NOT} SensorC)) -> X \textbf{NOT} CompletionLight)} \\
P15 & Light hold & While the product stays at C the light stays on. & \texttt{G((Mode \textbf{AND} CompletionLight \textbf{AND} X(Mode \textbf{AND} SensorC)) -> X CompletionLight)} \\
P16 & Latched motion & Once moving, the motor keeps running until the next station sensor is reached. & \texttt{G((Mode \textbf{AND} MotorForward \textbf{AND} X(Mode \textbf{AND} \textbf{NOT} SensorB \textbf{AND} \textbf{NOT} SensorC)) -> X MotorForward)} \\
P17 & Manual buttons ignored & Jog buttons have no effect in automatic mode. & \texttt{G((Mode \textbf{AND} \textbf{NOT} MotorForward \textbf{AND} X(Mode \textbf{AND} ForwardButton \textbf{AND} \textbf{NOT} (SensorA \textbf{AND} ButtonA) \textbf{AND} \textbf{NOT} (SensorB \textbf{AND} ButtonB))) -> X \textbf{NOT} MotorForward)} \\
P18 & Light only via ButtonC & The light never turns on without a prior ButtonC at C. & \texttt{G((Mode \textbf{AND} \textbf{NOT} CompletionLight \textbf{AND} X(Mode \textbf{AND} \textbf{NOT} (SensorC \textbf{AND} ButtonC))) -> X \textbf{NOT} CompletionLight)} \\
\bottomrule
\end{longtable}
}

\noindent\textbf{Assumptions and requirement gaps.}
{\small\begin{enumerate}
\item \textbf{P7}: the requirement does not say what happens to the light in manual mode; we assume it is OFF (manual mode is for maintenance jogging).
\item \textbf{Mode switch during motion}: not specified. We assume switching to manual immediately stops automatic motion (P6 applies), and switching back to automatic resets the sequence to "waiting at A".
\item \textbf{Out-of-sequence buttons} (e.g. ButtonB while at A) are assumed to be ignored; this is the only reading consistent with the station-by-station description.
\item \textbf{Stop at B/C}: the requirement implies the motor stops when the \emph{next} station sensor is reached; P10/P12 formalise that.
\item \textbf{Sensor dropout during motion}: sensors are assumed to be FALSE between stations (TC02 cycles 3 and 7).
\end{enumerate}}

\noindent\textbf{Validation.} reference implementation, simulation 6/6 cases PASS; mutation testing 2/2 mutants detected; 6 cases / 64 scan cycles.

\subsection{M06\_SlabThick --- Slab Thickness Alarm and Fault History}
\label{spec:M06_SlabThick}

\noindent \texttt{M06\_SlabThick}; 14 properties; source \texttt{specs/M06\_SlabThick\_spec.md}.

{\small
\begin{longtable}{@{}p{0.05\linewidth}p{0.13\linewidth}p{0.30\linewidth}>{\raggedright\arraybackslash}p{0.42\linewidth}@{}}
\toprule
\textbf{ID} & \textbf{Category} & \textbf{Requirement} & \textbf{LTL} \\
\midrule
\endhead
P1 & Alarm trigger & A valid fault activates the alarm in the same cycle. & \texttt{G(validFault -> THICKNESS\_ALARM)} \\
P2 & Minimum hold & A triggered alarm stays active for at least 500 ms. & \texttt{G((THICKNESS\_ALARM \textbf{AND} holdAcc < 500ms \textbf{AND} X THICKNESS\_SIGNAL\_VALID) -> X THICKNESS\_ALARM)} \\
P3 & Alarm hold & The alarm stays active while the fault persists. & \texttt{G((THICKNESS\_ALARM \textbf{AND} X validFault) -> X THICKNESS\_ALARM)} \\
P4 & Alarm release & After the fault clears and 500 ms have elapsed, the alarm is released. & \texttt{G((\textbf{NOT} THICKNESS\_FAULT \textbf{AND} holdAcc >= 500ms) -> \textbf{NOT} THICKNESS\_ALARM)} \\
P5 & Invalid signal & An invalid signal switches the alarm off immediately and clears its holding state. & \texttt{G(\textbf{NOT} THICKNESS\_SIGNAL\_VALID -> \textbf{NOT} THICKNESS\_ALARM)} \\
P6 & Counting & The counter increments exactly on each FALSE$\rightarrow$TRUE transition of a valid fault. & \texttt{G((\textbf{NOT} validFault \textbf{AND} X validFault) -> X Slab\_Fault\_Counter = Slab\_Fault\_Counter + 1)} \\
P7 & Counting & The counter does not change otherwise. & \texttt{G(\textbf{NOT} (\textbf{NOT} validFault \textbf{AND} X validFault) -> X Slab\_Fault\_Counter = Slab\_Fault\_Counter)} \\
P8 & History & On a new fault the timestamp is stored at index 0 and older entries shift up. & \texttt{G((\textbf{NOT} validFault \textbf{AND} X validFault) -> X(Slab\_Fault\_History[0] = FaultTimestamp \textbf{AND} Slab\_Fault\_History[k] = old Slab\_Fault\_History[k-1] for k = 1..4))} \\
P9 & History & Only the five most recent entries are retained. & implied by P8 (entry 4 is discarded) \\
P10 & No recording & No fault is recorded while the signal is invalid. & \texttt{G(\textbf{NOT} THICKNESS\_SIGNAL\_VALID -> (X Slab\_Fault\_Counter = Slab\_Fault\_Counter \textbf{AND} X Slab\_Fault\_History[0] = Slab\_Fault\_History[0]))} \\
P11 & Runtime & Runtime accumulates only while cutting, with a valid signal and no fault. & \texttt{G((CUTTING\_RUN \textbf{AND} THICKNESS\_SIGNAL\_VALID \textbf{AND} \textbf{NOT} THICKNESS\_FAULT) -> X Slab\_Runtime\_Accumulator = Slab\_Runtime\_Accumulator + CycleTime)} \\
P12 & Runtime & Otherwise the accumulated value is preserved. & \texttt{G(\textbf{NOT} (CUTTING\_RUN \textbf{AND} THICKNESS\_SIGNAL\_VALID \textbf{AND} \textbf{NOT} THICKNESS\_FAULT) -> X Slab\_Runtime\_Accumulator = Slab\_Runtime\_Accumulator)} \\
P13 & Init & At power-up the alarm is off and all counters, runtime and history entries are zero. & \texttt{\textbf{NOT} THICKNESS\_ALARM \textbf{AND} Slab\_Fault\_Counter = 0 \textbf{AND} Slab\_Runtime\_Accumulator = 0 \textbf{AND} FORALL k: Slab\_Fault\_History[k] = 0} (initial state) \\
P14 & Monotonic & The fault counter never decreases. & \texttt{G(X Slab\_Fault\_Counter >= Slab\_Fault\_Counter)} \\
\bottomrule
\end{longtable}
}

\noindent\textbf{Assumptions and requirement gaps.}
{\small\begin{enumerate}
\item \textbf{Units}: \texttt{CycleTime} in ms; \texttt{Slab\_Runtime\_Accumulator} is compared in \textbf{milliseconds} (toolchain convention).
\item \textbf{Hold restart}  Whether a fault recurring inside the 500 ms window restarts the hold is unspecified. The team adopted a \textbf{TOF} timer (the original window simply expires), so the tests drive it through \texttt{hold\_ms}. \emph{(flagged as ambiguous)}
\item \textbf{Validity returning with a fault already TRUE} is treated as a new fault edge, because the previous state was cleared when the signal became invalid (TC05).  \emph{(flagged as ambiguous)}
\item \textbf{History is read-only to callers}; the specification does not constrain external writes.
\end{enumerate}}

\noindent\textbf{Validation.} reference implementation, simulation 5/5 cases PASS; mutation testing 2/2 mutants detected; 5 cases / 32 scan cycles.

\subsection{M07\_HVAC --- Office HVAC Multi-mode Energy-Saving Control}
\label{spec:M07_HVAC}

\noindent \texttt{M07\_HVAC}; 13 properties; source \texttt{specs/M07\_HVAC\_spec.md}.

{\small
\begin{longtable}{@{}p{0.05\linewidth}p{0.13\linewidth}p{0.30\linewidth}>{\raggedright\arraybackslash}p{0.42\linewidth}@{}}
\toprule
\textbf{ID} & \textbf{Category} & \textbf{Requirement} & \textbf{LTL} \\
\midrule
\endhead
P1 & Fan condition & The fan never runs when unoccupied. & \texttt{G(\textbf{NOT} OCCUPANCY\_DETECTED -> \textbf{NOT} SUPPLY\_FAN\_RUN)} \\
P2 & Night mode & The fan never runs in Night mode. & \texttt{G(night -> \textbf{NOT} SUPPLY\_FAN\_RUN)} \\
P3 & Delay & The fan runs only after the conditions have held continuously for 30 s. & \texttt{G(SUPPLY\_FAN\_RUN -> delayAcc >= 30s)} \\
P4 & Reset & Loss of occupancy or selecting Night mode resets the timing. & \texttt{G((\textbf{NOT} OCCUPANCY\_DETECTED \textbf{OR} night) -> X delayAcc <= CycleTime)} \\
P5 & Liveness & While the conditions hold for 30 s the fan starts. & \texttt{G((OCCUPANCY\_DETECTED \textbf{AND} \textbf{NOT} night \textbf{AND} delayAcc >= 30s) -> SUPPLY\_FAN\_RUN)} \\
P6 & Auto occupied & Auto mode, occupied: 0 \textdegree{}C. & \texttt{G((HVAC\_MODE\_SELECTOR = 0 \textbf{AND} OCCUPANCY\_DETECTED) -> TEMP\_SETPOINT\_ADJUST = 0)} \\
P7 & Auto unoccupied & Auto mode, unoccupied: +3 \textdegree{}C in energy-saving mode, otherwise +2 \textdegree{}C. & \texttt{G((HVAC\_MODE\_SELECTOR = 0 \textbf{AND} \textbf{NOT} OCCUPANCY\_DETECTED) -> (TEMP\_SETPOINT\_ADJUST = 3 <-> ENERGY\_SAVING\_MODE) \textbf{AND} (TEMP\_SETPOINT\_ADJUST = 2 <-> \textbf{NOT} ENERGY\_SAVING\_MODE))} \\
P8 & Comfort & Comfort mode: 0 \textdegree{}C occupied, +1 \textdegree{}C unoccupied. & \texttt{G(HVAC\_MODE\_SELECTOR = 1 -> (TEMP\_SETPOINT\_ADJUST = 0 <-> OCCUPANCY\_DETECTED) \textbf{AND} (TEMP\_SETPOINT\_ADJUST = 1 <-> \textbf{NOT} OCCUPANCY\_DETECTED))} \\
P9 & Economy & Economy mode: 0 \textdegree{}C occupied, +4 \textdegree{}C unoccupied. & \texttt{G(HVAC\_MODE\_SELECTOR = 2 -> (TEMP\_SETPOINT\_ADJUST = 0 <-> OCCUPANCY\_DETECTED) \textbf{AND} (TEMP\_SETPOINT\_ADJUST = 4 <-> \textbf{NOT} OCCUPANCY\_DETECTED))} \\
P10 & Night adjustment & Night mode: +4 \textdegree{}C regardless of occupancy. & \texttt{G(night -> TEMP\_SETPOINT\_ADJUST = 4)} \\
P11 & Energy-saving scope & The energy-saving selection affects Auto mode only. & \texttt{G(HVAC\_MODE\_SELECTOR <> 0 -> (TEMP\_SETPOINT\_ADJUST independent of ENERGY\_SAVING\_MODE))}, i.e. P8--P10 do not mention it \\
P12 & Range & The adjustment always stays within $\pm$4 \textdegree{}C. & \texttt{G(TEMP\_SETPOINT\_ADJUST >= -4 \textbf{AND} TEMP\_SETPOINT\_ADJUST <= 4)} \\
P13 & Init & At power-up the fan is OFF and the adjustment 0. & \texttt{\textbf{NOT} SUPPLY\_FAN\_RUN \textbf{AND} TEMP\_SETPOINT\_ADJUST = 0} (initial state) \\
\bottomrule
\end{longtable}
}

\noindent\textbf{Assumptions and requirement gaps.}
{\small\begin{enumerate}
\item \textbf{The adjustment is not delayed}: only the fan has a 30 s delay, so \texttt{TEMP\_SETPOINT\_ADJUST} follows the current inputs immediately (TC01 step 1).  \emph{(flagged as ambiguous)}
\item \textbf{Invalid selector values} are excluded by the requirement ("callers supply a valid value") and are not tested.
\item \textbf{Initial outputs} are OFF/0 until the first calculation, as stated in requirement 6.
\end{enumerate}}

\noindent\textbf{Validation.} reference implementation, simulation 3/3 cases PASS; mutation testing 2/2 mutants detected; 3 cases / 24 scan cycles.

\subsection{M08\_IPv4 --- IPv4 Parsing with Format and Range Checks}
\label{spec:M08_IPv4}

\noindent \texttt{M08\_IPv4}; 9 properties; source \texttt{specs/M08\_IPv4\_spec.md}.

{\small
\begin{longtable}{@{}p{0.05\linewidth}p{0.13\linewidth}p{0.30\linewidth}>{\raggedright\arraybackslash}p{0.42\linewidth}@{}}
\toprule
\textbf{ID} & \textbf{Category} & \textbf{Requirement} & \textbf{LTL} \\
\midrule
\endhead
P1 & Correct parse & A well-formed address produces the value in network byte order. & \texttt{G(wellFormed(ipText) -> ipValue = s1*2\textasciicircum{}24 + s2*2\textasciicircum{}16 + s3*2\textasciicircum{}8 + s4)} \\
P2 & Byte order & The first segment occupies the most significant 8 bits. & contained in P1 \\
P3 & Malformed & Any malformed input produces 0. & \texttt{G(\textbf{NOT} wellFormed(ipText) -> ipValue = 0)} \\
P4 & Segment count & Fewer or more than four segments produce 0. & \texttt{G(segmentCount(ipText) <> 4 -> ipValue = 0)} \\
P5 & Empty segment & An empty segment produces 0. & \texttt{G(hasEmptySegment(ipText) -> ipValue = 0)} \\
P6 & Non-digit & Any non-digit, non-period character produces 0. & \texttt{G(hasNonDigit(ipText) -> ipValue = 0)} \\
P7 & Range & A segment above 255 produces 0. & \texttt{G(EXISTS k: segment\_k > 255 -> ipValue = 0)} \\
P8 & Zero address & The valid address 0.0.0.0 also produces 0. & \texttt{G(ipText = "0.0.0.0" -> ipValue = 0)} \\
P9 & Determinism & The output depends only on the current input (no internal state). & \texttt{G((ipText = X ipText) -> (ipValue = X ipValue))} \\
\bottomrule
\end{longtable}
}

\noindent\textbf{Assumptions and requirement gaps.}
{\small\begin{enumerate}
\item \textbf{Leading zeros} ("010.001.000.009") are accepted as decimal, since the requirement only demands decimal digits and a value within 0..255.  \emph{(flagged as ambiguous)}
\item \textbf{Whitespace} is not a digit, so " 1.2.3.4" and "1.2.3.4 " are malformed.
\item \textbf{Segment length}: a segment of more than three digits always exceeds 255 and is therefore rejected by P7.
\item \textbf{P8 is intentionally ambiguous} in the requirement: a zero output cannot distinguish "valid 0.0.0.0" from "invalid"; this is explicitly allowed.
\end{enumerate}}

\noindent\textbf{Validation.} reference implementation, simulation 4/4 cases PASS; mutation testing 2/2 mutants detected; 4 cases / 23 scan cycles.

% 20 tasks, 264 properties, 88 test cases.
